\section{Capítulo~\ref{sec:motorlearning}. Aprendizaje motor}

La regla de mínimos cuadrados y la ventaja de un cable de error propio son teoremas; la estructura del circuito con cable de error, el debilitamiento por coincidencia, el ajuste de la traza al retardo, el efecto posterior y la recuperación espontánea de la habilidad están medidos; que todo el circuito funcione como aprendizaje supervisado y construya un modelo interno es la hipótesis principal; los números del modelo de dos procesos son cálculos con el modelo del libro.

{\footnotesize
\setlength{\tabcolsep}{3pt}\renewcommand{\arraystretch}{1.15}
\begin{longtable}{>{\raggedright}p{0.5cm}>{\raggedright}p{4.0cm}>{\raggedright}p{5.6cm}>{\raggedright}p{2.3cm}>{\raggedright\arraybackslash}p{1.7cm}}
\toprule
N.º & Afirmación & Experimento y qué se midió & Fuente & Estado \\
\midrule\endfirsthead
\toprule
N.º & Afirmación & Experimento y qué se midió & Fuente & Estado \\
\midrule\endhead
1 & La regla~\eqref{eq:lms} sigue en promedio el gradiente del cuadrado del error & Resultado de la teoría de los filtros adaptativos: una corrección proporcional a la entrada y al error de salida es un descenso estocástico por el gradiente del error cuadrático medio. & \cite{WidrowHoff1960} & teoría \\
2 & Con un cable de error por salida, la velocidad no depende del número de salidas; con un solo número común, cae con el número de neuronas ruidosas (proposición~\ref{prop:teachers}) & Curvas de aprendizaje de una red lineal: el aprendizaje por perturbaciones aleatorias de las neuronas es más lento que el gradiente exacto tantas veces como neuronas perturbadas haya. & \cite{Werfel2005} & teoría \\
3 & El circuito con cable de error funciona como aprendizaje supervisado (proposición~\ref{prop:errorwire}) & Teorías deducidas de la estructura del circuito. Los experimentos de las filas 7--10 concuerdan con ellas; que todo el circuito funcione exactamente así no está demostrado. & \cite{Marr1969, Albus1971} & hipótesis \\
4 & Capa de expansión: unos $7\cdot10^{10}$ pequeñas neuronas \nt{GLUT}, más que en todo el resto del cerebro & Recuento de núcleos en una suspensión de tejido disuelto: en el cerebelo humano hay $69\cdot10^9$ neuronas de las $86\cdot10^9$ de todo el cerebro. Estereología: $101\cdot10^9$ neuronas pequeñas en el cerebelo de hombres de edad avanzada. & \cite{Azevedo2009, Andersen1992cereb} & medido \\
5 & Elevar la dimensión de la entrada vuelve lineal la dependencia deseada & Teorema sobre el número de particiones: una suma ponderada de $n$ entradas con umbral realiza un etiquetado arbitrario de unos $2n$ vectores en posición general. Que el cerebelo aproveche justamente esto es parte de la hipótesis de la fila~3. & \cite{Cover1965} & teoría \\
6 & Unas $3\cdot10^7$ neuronas \nt{GABA} de salida, cada una con del orden de $10^5$ entradas & Estereología del cerebelo humano: $30.5\cdot10^6$ neuronas de salida. Microscopía electrónica, rata: unas $1.75\cdot10^5$ entradas de la capa de expansión por neurona de salida. El neurotransmisor inhibidor de las neuronas de salida, en la revisión. & \cite{Andersen1992cereb, NapperHarvey1988, Ito2001} & medido \\
7 & Exactamente un cable de error por neurona de salida, alrededor de una vez por segundo, cada vez una descarga potente & Revisión de registros: cada neurona de salida recibe una única entrada trepadora \nt{GLUT}, que provoca una espiga compleja con una frecuencia del orden de $1$~Hz. & \cite{Ito2001} & medido \\
8 & La probabilidad de la descarga depende de la dirección del error del movimiento & Mono, región oculomotora del cerebelo, adaptación de sacadas: la dirección del error visual fija la probabilidad de la espiga compleja, y su magnitud, su momento; la descarga cambia las espigas simples en el intento siguiente y desplaza el movimiento a lo largo del error preferido de la célula. & \cite{Herzfeld2018} & medido \\
9 & La coincidencia de la descarga de error con una entrada activa debilita la entrada ($-\eta e_i x_j$) & Revisión de experimentos: la estimulación conjunta de una entrada de la capa de expansión y del cable de error provoca un debilitamiento duradero de esa entrada; la estimulación de la entrada sola, no. & \cite{Ito2001} & medido \\
10 & La duración de la traza está ajustada al retardo del error: con un retardo de $100\ms$ se debilita sobre todo la entrada activa $100\ms$ antes & Cortes y ratones vivos: en la región responsable del aprendizaje oculomotor, el debilitamiento está finamente ajustado a un intervalo de unos $120\ms$ entre la entrada y el error, lo que tarda la señal de error en esa tarea; en otra región, distintas células se ajustan a distintos intervalos. & \cite{Suvrathan2016} & medido \\
11 & El circuito es un filtro adaptativo~\eqref{eq:adaptivefilter} que aprende un modelo interno del cuerpo & Modelo deducido de la estructura del circuito. No hay un experimento directo en el que el filtro aprendido se haya medido como función de transferencia del cuerpo. & \cite{Fujita1982} & hipótesis \\
12 & La predicción resta las consecuencias de los propios movimientos; por eso uno no puede hacerse cosquillas & Personas: el tacto propio resulta menos cosquilloso que el mismo tacto ajeno; en el escáner, la corteza somatosensorial responde más débilmente al tacto propio, y el cerebelo está menos activo cuando el movimiento toca la piel. La explicación por resta de la predicción es una hipótesis. & \cite{Blakemore1998} & hipótesis \\
13 & Con una fuerza externa el fallo disminuye, y al retirarla la mano falla hacia el otro lado & Personas que mueven el mango de un robot en un campo de fuerzas: las trayectorias vuelven a ser rectas; al retirar el campo, el efecto posterior es casi la imagen especular de las primeras distorsiones, y se transfiere también a zonas del espacio no entrenadas. & \cite{ShadmehrMussaIvaldi1994} & medido \\
14 & Aprenden dos procesos~\eqref{eq:tworate}; tras la perturbación inversa, la habilidad vuelve sola & Personas, campo de fuerzas, luego un breve campo inverso e intentos con el error fijado: la corrección vuelve sola a la anterior; el modelo de dos procesos lo reproduce, el de uno, no. & \cite{Smith2006} & medido \\
15 & Números de la fig.~\ref{fig:tworate}: $0.84$ (lento, $0.63$) en $200$ movimientos; $6$ inversos; rápido, $-0.53$; lento, $0.46$; regreso hasta $0.37$ en $40$ movimientos & Cálculo con \texttt{models/\allowbreak motor\_\allowbreak learning.py}, $A_f=0.92$, $B_f=0.1$, $A_s=0.996$, $B_s=0.02$: $0.840$ ($0.634$ y $0.205$), $6$ movimientos, $-0.531$ y $0.457$, pico de $0.371$ en el movimiento $40$. & \texttt{models/\allowbreak motor\_\allowbreak learning.py} & modelo \\
16 & Hacen falta dos procesos porque el cuerpo cambia a distintas velocidades & Teoría: la estimación óptima con perturbaciones de varias duraciones da varios filtros con distintas constantes de tiempo y explica la recuperación espontánea y el reaprendizaje acelerado. & \cite{Kording2007} & hipótesis \\
\bottomrule
\end{longtable}}
